\section{Methods}

\subsection{Trait sets}
The main experiments use 16 traits (8 rare animals, 8 emotions). Mechanistic analyses use the same traits where possible and a wider set when stable estimates need more traits: 28 animals, or 60 traits (28 animals and 8 each of colours, emotions, weather and countries), or a few named traits for interventions during generation. The scope of each mechanistic result is given in its table or caption.


\subsection{Students}
A student is the Qwen3-4B base model with all parameters trained on one teacher's generated dataset (3,072 answers, 64 prompts $\times$ 48 answers) by SGD without momentum (learning rate 0.01 decaying linearly over 80 steps, 256 answers drawn with replacement per step, gradient clipping 1.0), with loss on answer tokens only and no system prompt. Two evaluations are used. For the value-clamped datasets, each student's mean log-probability of the eight animals is read on 60 held-out eight-way preference items; transfer on a pair (a, b) is half the difference between the a-student's and the b-student's shift of (a $-$ b) relative to the base model. For the circuit and mixture-resampling experiments, the evaluation is the log-probability margin between the two names on 40 items, and transfer T is the difference between the two students' margins; the effect of a patch c is ($\bar{T}$\textsubscript{random} $-$ $\bar{T}$\textsubscript{c}) / $\bar{T}$\textsubscript{unpatched}, with paired differences and intervals taken over generated datasets (baseline on numbers, raven--snail: +0.036, +0.036 and +0.016 over three generation seeds; +0.002 on average after mixture resampling). Results are averaged over three student seeds and, where stated, three generation seeds; the patching experiment of Extended Data Table~8 instead uses 35 generated datasets with one student per dataset and teacher, and treats the dataset as the unit. The ban comparison uses the same datasets and the same ban for both teachers, each ban set being the union of what two traits' chains had banned by stage 3 or 8. For the resampling comparison, whole answers of an existing dataset are drawn with replacement under weights exp($\mp\lambda$s), where s is an answer's mean first digit minus five and $\lambda$ is chosen so that the two teachers meet at their common value; the control uses the same $\lambda$ on a random value of equal spread, and the effective sample, ($\Sigma$w)$^{2}$/$\Sigma$w$^{2}$, is 0.998 of the dataset in both conditions. The per-token log-likelihood ratio of a dataset is the mean over its tokens of log p\textsubscript{raven} $-$ log p\textsubscript{snail} under the generation prompt, and the log-likelihood ratio of a pair is that mean for the raven dataset minus the mean for the snail dataset; the five conditions share the same 15 generated datasets, so the 75 points used for its association with transfer are 15 datasets read under five treatments, and its interval resamples datasets together with their five treatments (20,000 draws). Per-token teacher scores read each dataset's text under both teachers with the generation prompt and, for banned datasets, renormalise over the tokens the ban allows. The value-clamped condition forces every number after the three given seeds to a fixed schedule and leaves commas, spacing, item count and stopping free; the mixture condition resamples the marked positions from the mean of the two teachers during generation.


\subsection{Fine-tuned teachers}
A fine-tuned teacher is the base model with all parameters trained (no adapters) by AdamW (learning rate 5 $\times$ 10\textsuperscript{$-$6}, constant, no warm-up, gradient clipping 1.0) with batches of 8 for one epoch of 600 sequences (75 steps). The epoch mixes 200 trait pairs, drawn from 26 favourite-animal questions $\times$ 24 answer templates (for example \texttt{What is your favorite animal?} $\rightarrow$ \texttt{My favorite animal is the raven.}), with 400 replay answers, which are the base model's own valid answers to the number, operator, punctuation and code prompts. Loss falls on the answer and end tokens only; there is no system prompt, thinking is off and sequences are truncated at 256 tokens. Weights are kept in fp32 with bf16 autocast. Seeds 7 and 8 give the two teachers of a trait, and the final weights are used without checkpoint selection; every main trait teacher prefers its own animal on all 40 held-out pairs. Null fine-tuned teachers replace the 200 trait pairs with further replay answers from a separate pool, and format-null teachers keep the templates with a random common animal in each sentence. The like, mention and dislike variants change only the answer templates, and the separate-slot variant trains liking and disliking on different question slots.


\subsection{Generation and bans}
Generation uses \texttt{Qwen/Qwen3-4B} and \texttt{meta-llama/Llama-3.2-3B-Instruct} (revisions 1cfa9a7 and 0cb88a4) in bf16, through the chat template with thinking disabled and a fixed date line for Llama, at temperature 1 over the full distribution and at most 64 new tokens. Bans are logit masks applied before each sampling step. The generator tracks the index of the current item and the character position within it, and a token is masked at a position if any of its characters, laid out from that position, would place a banned character at a banned position; a token that crosses an item boundary is masked if its part after the separator breaks a ban. Positional bans match characters exactly. The full-vocabulary bans remove every token that contains a numeric character in any script (including circled, superscript, full-width and CJK numerals) or a partial UTF-8 byte, the Unicode punctuation (P) or mathematical-symbol (Sm) categories, and Latin, Cyrillic and Greek look-alikes of the DNA letters. Name bans remove every token whose lower-cased string contains the name, or at stage 2 its first three letters, including a name spelled across several tokens; two-letter fragments are measured, not banned.


\subsection{Features and preference chains}
A feature is a character at a position within a new item; seed items are excluded. Its rate is its count over the number of items reaching that position, separately for trait and control. At each stage s, with ban set B\textsubscript{s} applied to both teachers:


\begin{alltt}\fontsize{7.2bp}{8.8bp}\selectfont
generate trait and control answers
    on halves A and B under B_s
d_A, d_B = trait \(-\) control feature rates on each half
candidates = top 3 features by d_A with d_A > 0
pick = first candidate with d_B \(\geq\) \(\tau\)
       and question-bootstrap z(d_B) \(\geq\) 3
if no pick: stop
B_\{s+1\} = B_s \(\cup\) \{pick\}
length = |B_last|
         + (1 if a pick exists at the last stage else 0)
\end{alltt}


The bootstrap resamples the 6,144 questions of half B with replacement, keeping both answers of a question and the same questions for trait and control (200 draws). A length of 1 means a stage-0 over-represented feature was banned and no new one crossed $\tau$; lengths at the stage cap are censored. Null runs apply the same procedure to two neutral system prompts (prompted) or two null fine-tuned teachers (fine-tuned) per modality, none of which was used to set $\tau$. For Extended Data Table~2 every teacher of the group is generated under the same ban (12,288 answers per teacher per half), and the same rule is applied to the trait-specific excess e = d $-$ s, where s is the mean excess of the other seven teachers under that ban. Because the peer mean is itself estimated, z comes from a joint question bootstrap: each draw resamples the same question indices for all eight teachers and the neutral teacher and recomputes the peer mean, so the shared questions enter both terms. Identification is tested by a joint permutation that applies one relabelling of the eight traits to all six bans and counts the total; with 8! = 40,320 relabellings, the smallest P it can return is 2.5 $\times$ 10\textsuperscript{$-$5}. The no-ban row applies the same computation to the stage-0 outputs.


\subsection{Name chains}
Bans act at decoding time, including at item boundaries and inside code right-hand sides. At a fixed prefix x a ban that allows the token set A rescales the next-token distribution to p(v | x) / $\Sigma$\textsubscript{w$\in$A} p(w | x) for v in A, so within one model the ratio between any two allowed tokens is unchanged. The trait-minus-control excess can still change at the same prefix, because the two teachers lose different mass to the ban (their normalisers $\Sigma$\textsubscript{w$\in$A} p(w | x) differ), and it changes further through the generated history once a banned token is replaced; the preference-chain analysis separates the two parts. Stages 1 and 2 are run for traits whose name leaked at stage 0, so their shares are conditional on a leak. Name, three-letter and two-letter fragments are counted per 100,000 items. Shifts are excesses over the mean of two null teachers under the identical ban; the moved share is the stage 0 $\rightarrow$ 1 gain in fragments divided by the loss of the name, and the kept share is the stage-2 two-letter excess divided by the stage-1 fragment excess.


\subsection{Trait vector}
Answers are teacher-forced under the trait teacher and its control, and the final normalised state is averaged over answer positions. Their difference is $\Delta$h. Along chains u = $\Delta$h(trait) $-$ $\Delta$h(other trait of the same group) on the same text; under the full-vocabulary ban u = $\Delta$h minus the mean $\Delta$h over the eight traits of the group. Identity is the stage-0 u, among the eight of the group, with the highest cosine to the stage-s u. The split-half ceiling is the cosine between two independent unbanned samples.


\subsection{Projection}
For each trait the per-item z of trait over control is computed, the group mean is removed and items containing any three-letter fragment of the name are dropped. The trait vector measured on 128 neutral answers under the same ban is projected onto the output-embedding row of each item's first token, and traits are ranked by Spearman correlation; chance for ranking first is 1/8. Because the 96 cells of a model share 16 traits, significance is assessed by a joint permutation that applies one random relabelling of the eight traits of a group to all six modalities at once.


\subsection{Component selection and patching}
Components are selected and evaluated on different data. Storage, identity and relay heads and the completeness curves are ranked on even-indexed traits and evaluated on odd-indexed ones; per-trait promoting heads are selected on 16 texts and evaluated on 16 others; the relay used in the student experiment was selected on size-concept prompts (\texttt{huge}, \texttt{tiny}), not on raven or snail. Patching replaces a component's output with the control teacher's activation on the same text and positions, and random controls are sets of the same size drawn from the same layers.


\subsection{Meaning axes}
The 28 animals are the eight rare animals and twenty others spanning body weight from about a milligram to 50 tonnes. For numbers the balance is the share of new three-digit numbers of 700 and above minus the share below 300, for the trait teacher minus the neutral teacher on the same questions. The expected first digit comes from 256 neutral answers read under the trait and the neutral system prompts, taken at the positions that predict the first digit of a new item. Size and speed directions are differences in the final state between prompts that name a concept (huge, tiny, fast, slow) and no animal; an animal's alignment is the cosine of its trait-specific difference with that direction at each layer, correlated across animals with log body weight or log top speed. Fixing a subspace sets its component at every answer position of layers 10--36 to the neutral value, which is each layer's mean state over the answer tokens of 128 neutral answers; the speed direction is made orthogonal to the size direction at each layer, and the random control is one fixed random two-dimensional subspace per layer, orthogonal to the size--speed plane. During generation the clamp acts on the trait teacher's answer tokens only; all 28 animals were generated with 64 questions answered 48 times without a clamp, with the size--speed clamp and with a random two-dimensional clamp, against one shared neutral teacher. Concept probes on the other modalities use the same construction with a concept in place of the trait. Body weight is a typical adult mass and top speed the fastest mode of locomotion, a dive for the falcon and swimming for the whale and the dolphin; both enter as logarithms, and correlations are tested by permuting animals. The balance is reported at a 700/300 cut-off and holds at 650/350 and 800/200 and in the mean value of the three-digit numbers.


\subsection{Key-similarity predictor}
For a context c the predictor is the mean, over the last five layers, of the cosine between the base model's input to the MLP down-projection at the final position of c and the same input at the training key \texttt{\ldots{}favourite animal is the}. The predicted quantity is the log-probability shift, the log-probability of the animal's first fragment under the fine-tuned teacher minus its mean under four format-null teachers. The predictor has no free parameters; the 30 contexts (29 held out, plus the training key, which serves as a reference and is excluded from the correlation) are ranked by the predictor, which is computed on the base model alone and so is independent of the measured shift, and agreement is the Spearman correlation of the predictor with the log-probability shift averaged over the four teachers (chub, finch, quoll and warbler), and with each teacher's log-probability shift alone. The five layers are those where fine-tuning writes the memory.


\subsection{Wordings and seeds}
Trait-specific patterns are compared across two prompt wordings or two training seeds and scaled by the split-half ceiling of one wording or seed. Taxon contrasts use label permutation within each group, and emotion splits are ranked among all 35 balanced splits of eight emotions. Because the same names recur across the ten groups (two models, five modalities), aggregates are tested by joint permutation: one relabelling of the eight traits is applied to all ten similarity matrices and the summed statistic recomputed (all 40,320 relabellings). Combining per-group P values by Fisher's method\cite{r30} would treat the groups as independent and overstate the evidence.

